% =====================================================================
%  KnotBench -- NeurIPS / ICML submission skeleton
%  Track: Datasets & Benchmarks (default). See \maintrack switch below.
%  Every table is generated: python scripts/make_paper_tables.py (results/benchmark_v2/*.json).
%  \todo{} marks items blocked by gaps listed in the paper blueprint.
% =====================================================================
\documentclass{article}

% --- Style: drop the official file (e.g. neurips_2027.sty) next to this .tex.
%     For D&B use the track option the call specifies (e.g. [datasets]/[db]).
\IfFileExists{neurips_2027.sty}
  {\usepackage{neurips_2027}}               % anonymous submission by default
  {\usepackage[margin=1in]{geometry}\usepackage[numbers]{natbib}}  % fallback so the draft compiles

\usepackage[utf8]{inputenc}
\usepackage[T1]{fontenc}
\IfFileExists{lmodern.sty}{\usepackage{lmodern}}{}
\usepackage{hyperref}
\usepackage{url}
\usepackage{booktabs}
\usepackage{amsfonts,amsmath,amssymb,amsthm}
\usepackage{nicefrac}
\usepackage[expansion=false]{microtype}
\usepackage{xcolor}
\usepackage{graphicx}
\usepackage{multirow}
\usepackage{enumitem}
\usepackage{adjustbox}
\usepackage{longtable}

\newcommand{\todo}[1]{\textcolor{red}{[TODO: #1]}}
\newcommand{\flag}[1]{\textcolor{gray}{#1}}   % greyed, non-competing rows
\newtheorem{proposition}{Proposition}

% Set to true to build the main-track (methods) variant
\newif\ifmaintrack \maintrackfalse

\title{KnotBench: A Leakage-Free Benchmark for\\Machine Learning on Knot Diagrams}

\author{Anonymous Authors}  % keep anonymous for review

\begin{document}
\maketitle

% ---------------------------------------------------------------------
\begin{abstract}
Machine learning on knot diagrams is increasingly used to predict knot
invariants and to decide knot equivalence, yet reported gains are rarely
measured against exact simplification, classical invariants, or tree models
given the same inputs. We introduce \textsc{KnotBench}, built entirely from
KnotInfo, with knot-level splits by crossing number (train $c{\le}10$,
validate $c{=}11$, test on 2{,}176 knots with $c{=}12$), two presentations per
knot, size-matched pairs, hard geometric negatives, and controls that isolate
what learning adds. Four parameter-matched graph backbones (signed MPNN, GIN,
GPS, and a node-marking subgraph GNN) beat diagram-only tree ensembles on
unknotting number and smooth 4-genus, but none beats a two-parameter
$[c,w]$ regression on signature or Seifert genus at 250 training knots; that
gap closes at 636 training knots. On hard negatives drawn from distinct
presentations, trained encoders reach only 0.45--0.54 AUROC while an exact
hyperbolic certificate scores 1.000 at 1.3\,ms per pair. We reproduce the
natural-slope result of Davies et al.\ and document five evaluation pitfalls
with measured effect sizes. Data, code and a one-command protocol are released.
\end{abstract}

% ---------------------------------------------------------------------
\section{Introduction}\label{sec:intro}
% Para 1: why ML for knots (Davies et al. 2021). Para 2: evaluation shortcuts
% inflate results (Fig. 1). Para 3: what KnotBench does. Para 4: findings.
\todo{prose}

\begin{figure}[t]
  \centering
  \fbox{\parbox{0.9\linewidth}{\centering\vspace{1.0cm}
  Fig.~1 (teaser): verification AUROC on hard geometric negatives per slice --
  face histogram (0.23--0.32), best untrained (0.40--0.44), best trained
  (0.48--0.54), exterior-isometry certificate (1.000).\vspace{1.0cm}}}
  \caption{On distinct presentations with hard negatives, learned encoders stay
  near chance while an exact certificate is perfect and fast. Data:
  \texttt{results/benchmark\_v2/task1\_hard\_negatives.json}.}
  \label{fig:teaser}
\end{figure}

\paragraph{Contributions.}
\begin{enumerate}[leftmargin=*,itemsep=2pt]
  \item \textbf{Benchmark.} Knot-level splits by crossing number from KnotInfo;
        verification across two presentations with size-matched and hard
        geometric negatives; invariant prediction with a held-out validation
        slice and no test-set tuning.
  \item \textbf{Fair and oracle baselines.} Diagram-only tree ensembles see the
        same information as the graph models; oracle rows that read classical
        invariants ($\sigma$, $\det$, Alexander) are reported separately.
  \item \textbf{Matched backbones and controls.} Four backbones at
        $\approx$110k parameters; untrained and shuffled-label controls;
        paired, Holm-corrected tests for every claimed difference.
  \item \textbf{Pitfall taxonomy and a reproduction} of the Davies et al.\
        signature result inside the leakage-free split.
  \ifmaintrack
  \item \textbf{Method.} \todo{non-local architecture}
  \fi
\end{enumerate}
Code and data: \todo{anonymous link}.

% ---------------------------------------------------------------------
\section{Background and Related Work}\label{sec:related}
\paragraph{Knot diagrams and invariants.} PD codes, Reidemeister moves
$R_1$--$R_3$, flypes~\cite{menasco1993classification}; targets $\sigma$, $g_3$,
$g_4$, $u$.
\paragraph{Machine learning for knots.} Davies et al.~\cite{davies2021advancing}
related signature to cusp geometry, later proved as a natural-slope
inequality~\cite{davies2022signature}. \todo{Hughes; Jejjala et al.; Gukov et al.}
\paragraph{GNN expressivity.} 1-WL bound~\cite{xu2019powerful}; subgraph
GNNs~\cite{bevilacqua2022equivariant}; graph transformers~\cite{rampasek2022gps}.
\paragraph{Benchmark leakage.} \todo{}

% ---------------------------------------------------------------------
\section{Benchmark Design}\label{sec:benchmark}
\subsection{Data and splits}
All 2{,}978 knots up to 12 crossings and their labels come from
KnotInfo~\cite{knotinfo}. Invariant prediction trains on $c{\le}10$ (250
knots), selects checkpoints per target on $c{=}11$ (552) and tests on $c{=}12$
(2{,}176; $u(K)$ on the 1{,}628 with an exact value).
\subsection{Verification pairs}
Both presentations of a knot (PD code and braid closure) come from KnotInfo,
which avoids the Perko-pair renumbering at $10_{161}$--$10_{165}$. Diagrams are
$R_1/R_2$-reduced and relabelled. Hard negatives are the two nearest
neighbours in face-degree histogram space $[c,n_+,n_-,w,h_2,\dots,h_{10}]$.
Encoders train on all knots with $3{\le}c{\le}8$ plus alternating knots with
$c\in\{9,10\}$ (199 knots), knot-disjoint from the non-alternating test knots.
\subsection{Metrics and statistics}
AUROC and MAE, mean $\pm$ std over seeds $\{42,\dots,46\}$. Differences are
tested with paired knot-level bootstrap and Holm correction over all 52
comparisons (Appendix~\ref{app:stats}).

% ---------------------------------------------------------------------
\section{Models, Baselines and Controls}\label{sec:models}
\begin{proposition}[Confluent $R_1/R_2$ reduction]
Exhaustive $R_1$-curl and $R_2$-bigon removal terminates in $O(|V|)$ steps and
its result is independent of move order. \todo{state precisely; proof in
Appendix~\ref{app:proof}}
\end{proposition}
\begin{itemize}[leftmargin=*,itemsep=1pt]
  \item \textbf{Backbones} (hidden width chosen so each has $\approx$110k
        parameters; 4 layers; AdamW, 60 epochs, batch 32): signed MPNN,
        GINE~\cite{xu2019powerful}, GPS~\cite{rampasek2022gps}, DS-GNN with
        node marking~\cite{bevilacqua2022equivariant}.
  \item \textbf{Diagram-only trees}: RF and GBDT on the 13 diagram features above.
  \item \textbf{Oracle trees} (not competitors): diagram features plus
        $\sigma$, $|\sigma|/2$, $\det$, Alexander coefficients; $\sigma$ columns
        removed when $\sigma$ is the target.
  \item \textbf{Classical}: WL hash, $|\sigma|/2$ bound, $\deg\Delta/2$, exterior
        isometry certificate (ceiling).
\end{itemize}

% ---------------------------------------------------------------------
\section{Results}\label{sec:results}

\begin{table}[t]
\centering\small
\caption{Invariant prediction on unseen $c{=}12$ knots (train $c{\le}10$,
checkpoint on $c{=}11$), MAE, mean $\pm$ std over 5 seeds. Bold: best non-oracle
entry. Grey rows read classical invariants and are not competitors.}
\label{tab:task2}
\begin{adjustbox}{max width=\linewidth}
\begin{tabular}{l r cccccccc}
\toprule
 & & \multicolumn{2}{c}{Trivial} & \multicolumn{2}{c}{Diagram-only trees} & \multicolumn{4}{c}{Graph backbones (param.-matched)}\\
\cmidrule(lr){3-4}\cmidrule(lr){5-6}\cmidrule(lr){7-10}
Target & $N$ & Mean & Lin.\ $[c,w]$ & RF & GBDT & Signed MPNN & GIN & GPS & Subgraph (DS-GNN)\\
\midrule
$u(K)$ & 1628 & 0.489 & 0.471 & $0.432\pm0.003$ & $0.426\pm0.003$ & \textbf{$0.396\pm0.013$} & $0.408\pm0.017$ & $0.436\pm0.018$ & $0.401\pm0.008$\\
$g_4(K)$ & 2176 & 0.597 & 0.524 & $0.388\pm0.009$ & $0.385\pm0.004$ & $0.363\pm0.018$ & \textbf{$0.361\pm0.008$} & $0.366\pm0.011$ & $0.364\pm0.022$\\
$g_3(K)$ & 2176 & 0.855 & \textbf{0.619} & $0.744\pm0.003$ & $0.723\pm0.014$ & $0.729\pm0.022$ & $0.669\pm0.024$ & $0.747\pm0.038$ & $0.693\pm0.074$\\
$\sigma(K)$ & 2176 & 2.325 & \textbf{0.921} & $0.989\pm0.007$ & $0.974\pm0.011$ & $1.004\pm0.041$ & $0.947\pm0.020$ & $1.008\pm0.044$ & $0.974\pm0.048$\\
\midrule
\multicolumn{10}{l}{\emph{Oracle inputs (not competitors): trees given $\sigma$, $\det$, Alexander coefficients; classical formulas}}\\
\flag{$u(K)$} & & \multicolumn{8}{l}{\flag{RF $0.372\pm0.011$, GBDT $0.373\pm0.011$; closed form 0.848}}\\
\flag{$g_4(K)$} & & \multicolumn{8}{l}{\flag{RF $0.137\pm0.006$, GBDT $0.310\pm0.029$; closed form 0.271}}\\
\flag{$g_3(K)$} & & \multicolumn{8}{l}{\flag{RF $0.056\pm0.001$, GBDT $0.056\pm0.000$ (formula input); closed form 0.015}}\\
\flag{$\sigma(K)$} & & \multicolumn{8}{l}{\flag{RF $0.857\pm0.005$, GBDT $0.858\pm0.008$ ($\sigma$ columns removed)}}\\
\bottomrule
\end{tabular}

\end{adjustbox}
\end{table}

\begin{table}[t]
\centering\small
\caption{Verification across distinct presentations (AUROC, mean $\pm$ std over
5 seeds). Ceiling = exterior isometry certificate. Backbone columns show the
better of supervised and contrastive training; Untrained is the best of the four
untrained backbones.}
\label{tab:task1}
\begin{adjustbox}{max width=\linewidth}
% certificate: median 1.3 ms/pair; methods {'isometry': 1331, 'hfk_fallback': 5}
\begin{tabular}{l l c c c cccc}
\toprule
Slice & Neg. & Ceiling & Face hist. & Untrained & Signed MPNN & GIN & GPS & Subgraph (DS-GNN)\\
\midrule
$c\le10$ (50) & random & 1.000 & 0.937 & $0.675\pm0.017$ & $0.642\pm0.040$ & $0.840\pm0.097$ & $0.656\pm0.060$ & $0.656\pm0.048$\\
$c\le10$ (50) & hard & 1.000 & 0.321 & $0.428\pm0.024$ & $0.453\pm0.034$ & $0.483\pm0.040$ & $0.486\pm0.032$ & $0.483\pm0.035$\\
$c=11$ (185) & random & 1.000 & 0.867 & $0.677\pm0.005$ & $0.656\pm0.022$ & $0.802\pm0.074$ & $0.672\pm0.035$ & $0.693\pm0.025$\\
$c=11$ (185) & hard & 1.000 & 0.230 & $0.401\pm0.007$ & $0.477\pm0.034$ & $0.507\pm0.030$ & $0.493\pm0.022$ & $0.479\pm0.032$\\
$c=12$ (100) & random & 1.000 & 0.909 & $0.697\pm0.012$ & $0.750\pm0.045$ & $0.834\pm0.076$ & $0.742\pm0.039$ & $0.720\pm0.016$\\
$c=12$ (100) & hard & 1.000 & 0.267 & $0.441\pm0.015$ & $0.506\pm0.023$ & $0.543\pm0.059$ & $0.492\pm0.078$ & $0.481\pm0.013$\\
\bottomrule
\end{tabular}

\end{adjustbox}
\end{table}

\paragraph{Graph models beat diagram-only trees, not classical invariants.}
On $u(K)$ and $g_4(K)$ every backbone except GPS on $u$ beats diagram-only RF
(e.g.\ GIN $-0.035$ MAE on $u$, $-0.034$ on $g_4$; $p_\text{Holm}{=}0.010$).
Trees given $\sigma$ reach 0.137 on $g_4$, so the earlier ``RF beats GNN''
result came from oracle inputs.
\paragraph{Signature and Seifert genus are data-limited.} At 250 training
knots no diagram model beats linear $[c,w]$ on $\sigma$ (GIN $+0.007$, n.s.) or
$g_3$. With 636 training knots GIN reaches 0.718 on $\sigma$ and 0.472 on $g_3$
versus 0.919 and 0.613 for $[c,w]$ (Table~\ref{tab:curve}).
\paragraph{Learning adds little on hard negatives.} Training lifts the best
encoder from 0.40--0.44 to 0.48--0.54 AUROC (significant at $c{=}11,12$, not
at $c{\le}10$); all remain far below the certificate.
\paragraph{Cost.} The certificate costs 1.3\,ms per pair (median), comparable
to a single encoder forward pass (Table~\ref{tab:runtime}).
\paragraph{Reproducing Davies et al.} Hyperbolic invariants predict $\sigma$
far better than any diagram model (Table~\ref{tab:davies}).

\begin{table}[t]
\centering\small
\caption{Learning curve: test MAE on $c{=}12$ as the training set grows
(validation fixed at 30\% of $c{=}11$).}
\label{tab:curve}
\begin{adjustbox}{max width=\linewidth}
\begin{tabular}{l l ccc}
\toprule
Target & Model & c$\le$9 (85) & c$\le$10 (250) & c$\le$10+70\%c11 (636)\\
\midrule
$u(K)$ & Lin.\ $[c,w]$ & $0.540\pm0.000$ & $0.471\pm0.000$ & $0.445\pm0.000$\\
 & RF (diagram) & $0.477\pm0.009$ & $0.432\pm0.003$ & $0.423\pm0.002$\\
 & Signed MPNN & $0.450\pm0.010$ & $0.408\pm0.007$ & $0.404\pm0.014$\\
 & GIN & $0.469\pm0.035$ & $0.409\pm0.015$ & $0.388\pm0.007$\\
 & GPS & $0.436\pm0.018$ & $0.441\pm0.021$ & $0.427\pm0.007$\\
 & Subgraph (DS-GNN) & $0.448\pm0.016$ & $0.420\pm0.017$ & $0.427\pm0.007$\\
\midrule
$g_4(K)$ & Lin.\ $[c,w]$ & $0.575\pm0.000$ & $0.524\pm0.000$ & $0.514\pm0.000$\\
 & RF (diagram) & $0.441\pm0.008$ & $0.388\pm0.009$ & $0.385\pm0.003$\\
 & Signed MPNN & $0.437\pm0.003$ & $0.359\pm0.012$ & $0.330\pm0.009$\\
 & GIN & $0.419\pm0.010$ & $0.361\pm0.008$ & $0.314\pm0.009$\\
 & GPS & $0.432\pm0.022$ & $0.366\pm0.011$ & $0.347\pm0.012$\\
 & Subgraph (DS-GNN) & $0.431\pm0.012$ & $0.362\pm0.020$ & $0.342\pm0.019$\\
\midrule
$g_3(K)$ & Lin.\ $[c,w]$ & $0.594\pm0.000$ & $0.619\pm0.000$ & $0.613\pm0.000$\\
 & RF (diagram) & $1.036\pm0.008$ & $0.744\pm0.003$ & $0.613\pm0.002$\\
 & Signed MPNN & $0.963\pm0.070$ & $0.729\pm0.022$ & $0.574\pm0.022$\\
 & GIN & $0.912\pm0.077$ & $0.672\pm0.030$ & $0.472\pm0.028$\\
 & GPS & $1.024\pm0.031$ & $0.743\pm0.037$ & $0.586\pm0.012$\\
 & Subgraph (DS-GNN) & $0.978\pm0.055$ & $0.693\pm0.073$ & $0.564\pm0.024$\\
\midrule
$\sigma(K)$ & Lin.\ $[c,w]$ & $0.896\pm0.000$ & $0.921\pm0.000$ & $0.919\pm0.000$\\
 & RF (diagram) & $1.180\pm0.013$ & $0.989\pm0.007$ & $0.966\pm0.003$\\
 & Signed MPNN & $1.347\pm0.040$ & $1.013\pm0.042$ & $0.850\pm0.020$\\
 & GIN & $1.211\pm0.026$ & $0.950\pm0.021$ & $0.718\pm0.052$\\
 & GPS & $1.342\pm0.068$ & $1.007\pm0.044$ & $0.918\pm0.058$\\
 & Subgraph (DS-GNN) & $1.412\pm0.055$ & $0.971\pm0.051$ & $0.899\pm0.058$\\

\bottomrule
\end{tabular}

\end{adjustbox}
\end{table}

\begin{table}[t]
\centering\small
\caption{Signature from hyperbolic invariants of the complement (Davies et al.\
setting vs.\ the KnotBench crossing split); 5 seeds; ``exact'' rounds to the
nearest even integer.}
\label{tab:davies}
\begin{tabular}{l c c c c}
\toprule
 & \multicolumn{2}{c}{Random 80/20 (Davies et al.\ setting)} & \multicolumn{2}{c}{KnotBench $c\le10\to c=12$}\\
\cmidrule(lr){2-3}\cmidrule(lr){4-5}
Predictor of $\sigma(K)$ & MAE & Exact acc. & MAE & Exact acc.\\
\midrule
Mean & 2.292 & -- & 2.316 & --\\
Natural slope / 2 (formula) & 0.713 & 0.739 & 0.686 & 0.739\\
MLP (hyperbolic inv.) & $0.501\pm0.013$ & 0.862 & $0.596\pm0.011$ & 0.831\\
GBDT (hyperbolic inv.) & $0.424\pm0.003$ & 0.873 & $0.542\pm0.005$ & 0.808\\
RF (hyperbolic inv.) & $0.344\pm0.002$ & 0.868 & $0.476\pm0.002$ & 0.804\\
\bottomrule
\end{tabular}

\end{table}

\begin{table}[t]
\centering\small
\caption{Runtime on one CPU core.}
\label{tab:runtime}
\begin{tabular}{l r r}
\toprule
Method & Params & Time per diagram / pair\\
\midrule
Exterior isometry (certificate) & -- & 1.3 ms (median), 2.1 ms (p95)\\
WL hash & -- & 0.71 ms\\
Signed MPNN encoder & 123,207 & 0.57 ms\\
GIN encoder & 119,555 & 0.65 ms\\
GPS encoder & 131,535 & 0.54 ms\\
Subgraph (DS-GNN) encoder & 123,271 & 4.30 ms\\
\bottomrule
\end{tabular}

\end{table}

% ---------------------------------------------------------------------
\section{A Taxonomy of Pitfalls}\label{sec:pitfalls}
\begin{table}[h]
\centering\small
\caption{Pitfalls, how to detect them, and measured effect.}
\label{tab:pitfalls}
\begin{tabular}{p{3.2cm}p{4.2cm}p{4.3cm}}
\toprule
Pitfall & Detection / fix & Measured effect \\
\midrule
Traversal-order leakage & Cyclic relabel + orientation reversal; face-degree features & Chord AUROC $0.918\to0.401$ \\
Diagram-size imbalance & Match step schedules for positives/negatives & Spurious $0.000\to0.512$ AUROC \\
Encoder as graph hash & Untrained and shuffled-label controls; hard negatives & Untrained $0.95$ on easy pairs, $0.40$ on hard \\
Oracle / label leakage & Audit features against label definitions & $\sigma$ MAE $0.010$ (leak); $g_4$ $0.388\to0.137$ with $\sigma$ \\
Multi-task checkpointing & Per-target selection on validation & Joint vs.\ per-target selection \\
\bottomrule
\end{tabular}
\end{table}

% ---------------------------------------------------------------------
\section{Relevance to Scientific Practice}\label{sec:impact}
\textsc{KnotBench} answers a practical question for working scientists: when
should a learned model replace or complement an exact method? For knot
equivalence it should not: an exterior-isometry certificate is correct on every
test pair at 1.3\,ms per pair, while trained encoders stay near chance on hard
negatives. For invariants without fast algorithms---unknotting number and
smooth 4-genus---all four graph backbones carry significant signal beyond
diagram-only baselines, which makes them useful for ranking open cases (548 of
2{,}176 twelve-crossing knots lack an exact $u(K)$ in KnotInfo) but not for
settling them. The same protocol applies to knotted proteins, circular DNA and
polymer simulations, where knot identification runs at database scale: screen
with exact invariants, and adopt a learned model only if it beats them on accuracy
and cost under a knot-disjoint split. The oracle-input pitfall
($g_4$ MAE $0.388\to0.137$ once $\sigma$, which helps define the labels, is an input)
is a general hazard for scientific benchmarks whose labels come from bounds. For physical tangles (cables, ropes, threads) the same split applies:
learn perception, compute topology. A companion photo benchmark with the
same splits and metrics (crossing precision/recall, over/under accuracy, topology
accuracy with abstention, cable-level confidence intervals) is released with
synthetic renders. A skeleton-based parser recovers 16\% of crossings there; a
knot-aware NanoTransformer that predicts crossings, over-strand axes and cable masks
(instead of a knot label) and hands them to exact topology recovers
90\% of crossings with 100\% over/under
accuracy on matched crossings and answers 52\% of images with no wrong
answer (31\% on two unseen background styles), abstaining elsewhere.
These are synthetic results; real photographs are untested.

\section{Reproducibility Statement}\label{sec:repro}
All headline numbers are produced on CPU with seeds $\{42,\dots,46\}$ by one
script and written to JSON, from which every table is generated. Repeated runs,
different hash seeds and thread counts give bit-identical test predictions
(identical SHA-256 over all 2{,}176 predictions). The full suite takes about
2 CPU-hours with no GPU. One earlier value was platform-dependent: the learned
gain in the $c{\le}8$ protocol was $+0.008$ AUROC on Apple MPS with unseeded
initialisation and is $+0.027$ [$+0.013$, $+0.042$] on seeded CPU runs; we
report the latter. We release pinned dependencies (SnapPy, spherogram,
scikit-learn, database\_knotinfo), data with checksums, and record the commit,
package versions and data hash in every results file. \todo{confirm once the
lockfile, Dockerfile and DOI exist; add cross-platform tolerance}

% ---------------------------------------------------------------------
\section{Limitations}\label{sec:limits}
Crossing numbers $\le 12$; synthetic perturbations; KnotInfo labels for $u$ and
$g_4$ exist where classical bounds meet, which biases the solved subset; $R_3$
is not canonicalized; 5 of 1{,}336 certificate pairs are non-hyperbolic and use
knot Floer homology equality instead; the learned gain in the $c{\le}8$ protocol
(+0.027 AUROC on CPU) differs from the earlier MPS run (+0.008).

\paragraph{Topology is not effort.} \textsc{KnotBench} scores combinatorial
quantities: crossing number, unknotting number and equivalence. How hard a
physical cable, rope or thread is to undo also depends on its bending
stiffness, friction, tension and how tightly it is pulled (ropelength). A
diagram with fewer crossings can take more work to untie, and two tangles of the
same knot type can differ greatly in effort. Our untangling steps
(Reidemeister~I/II moves) and unknotting numbers are therefore lower bounds on
the number of manipulations, not predictions of force or time; physical claims
need mechanical models or experiments.

\section{Conclusion}\label{sec:conclusion}
\todo{two sentences + how to submit results to the benchmark}

% ---------------------------------------------------------------------
\bibliographystyle{plainnat}
\bibliography{references}

\appendix
\section{Proof of Confluent $R_1/R_2$ Reduction}\label{app:proof}
\todo{from docs/reidemeister\_invariance\_theory.md}
\section{Statistical Tests}\label{app:stats}
Paired knot-level bootstrap; $\Delta$ is MAE(A)$-$MAE(B) for Task 2 and
AUROC(A)$-$AUROC(B) for Task 1, computed on seed-averaged predictions; Holm correction over the full family. Resampling floors $p$ at $1/\text{reps}$.
{\footnotesize\begin{longtable}{p{8.2cm} r c r r}
\toprule
Comparison (A vs B) & $\Delta$ & 95\% CI & $p$ & $p_\text{Holm}$\\
\midrule
\endhead
gain / canonical\_r1r2 vs random\_canon & +0.027 & [+0.013, +0.042] & 0.0005 & 0.010\\
gain / canonical\_r1r2 vs shuffled\_canon & +0.029 & [+0.015, +0.045] & 0.0005 & 0.010\\
gain / canonical\_r1r2 vs fixed\_chord\_canon & +0.058 & [+0.038, +0.079] & 0.0005 & 0.010\\
task1 / c$\le$10\_hard\_geometric / gin\_contrastive vs gin\_untrained & +0.072 & [-0.029, +0.167] & 0.1680 & 1.000\\
task1 / c$\le$10\_hard\_geometric / gin\_contrastive vs face\_histogram & +0.175 & [+0.073, +0.270] & 0.0005 & 0.010\\
task1 / c=11\_hard\_geometric / gin\_supervised vs gin\_untrained & +0.106 & [+0.074, +0.139] & 0.0005 & 0.010\\
task1 / c=11\_hard\_geometric / gin\_supervised vs face\_histogram & +0.268 & [+0.224, +0.311] & 0.0005 & 0.010\\
task1 / c=12\_hard\_geometric / gin\_contrastive vs gin\_untrained & +0.135 & [+0.069, +0.202] & 0.0005 & 0.010\\
task1 / c=12\_hard\_geometric / gin\_contrastive vs face\_histogram & +0.309 & [+0.237, +0.380] & 0.0005 & 0.010\\
task2 / unknotting\_number / gnn\_mpnn vs rf\_diagram & -0.041 & [-0.057, -0.025] & 0.0002 & 0.010\\
task2 / unknotting\_number / gnn\_gin vs rf\_diagram & -0.035 & [-0.051, -0.019] & 0.0002 & 0.010\\
task2 / unknotting\_number / gnn\_gps vs rf\_diagram & -0.006 & [-0.022, +0.011] & 0.5028 & 1.000\\
task2 / unknotting\_number / gnn\_subgraph vs rf\_diagram & -0.038 & [-0.054, -0.021] & 0.0002 & 0.010\\
task2 / unknotting\_number / gnn\_mpnn vs linear\_cw & -0.081 & [-0.094, -0.069] & 0.0002 & 0.010\\
task2 / unknotting\_number / gnn\_gin vs linear\_cw & -0.075 & [-0.088, -0.062] & 0.0002 & 0.010\\
task2 / unknotting\_number / gnn\_gps vs linear\_cw & -0.046 & [-0.061, -0.031] & 0.0002 & 0.010\\
task2 / unknotting\_number / gnn\_subgraph vs linear\_cw & -0.078 & [-0.091, -0.065] & 0.0002 & 0.010\\
task2 / unknotting\_number / rf\_oracle vs rf\_diagram & -0.060 & [-0.071, -0.050] & 0.0002 & 0.010\\
task2 / unknotting\_number / rf\_diagram vs linear\_cw & -0.040 & [-0.057, -0.023] & 0.0002 & 0.010\\
task2 / smooth\_four\_genus / gnn\_mpnn vs rf\_diagram & -0.037 & [-0.051, -0.023] & 0.0002 & 0.010\\
task2 / smooth\_four\_genus / gnn\_gin vs rf\_diagram & -0.034 & [-0.047, -0.020] & 0.0002 & 0.010\\
task2 / smooth\_four\_genus / gnn\_gps vs rf\_diagram & -0.027 & [-0.040, -0.013] & 0.0002 & 0.010\\
task2 / smooth\_four\_genus / gnn\_subgraph vs rf\_diagram & -0.042 & [-0.055, -0.028] & 0.0002 & 0.010\\
task2 / smooth\_four\_genus / gnn\_mpnn vs linear\_cw & -0.175 & [-0.188, -0.162] & 0.0002 & 0.010\\
task2 / smooth\_four\_genus / gnn\_gin vs linear\_cw & -0.171 & [-0.186, -0.156] & 0.0002 & 0.010\\
task2 / smooth\_four\_genus / gnn\_gps vs linear\_cw & -0.165 & [-0.178, -0.151] & 0.0002 & 0.010\\
task2 / smooth\_four\_genus / gnn\_subgraph vs linear\_cw & -0.179 & [-0.193, -0.166] & 0.0002 & 0.010\\
task2 / smooth\_four\_genus / gnn\_subgraph vs gnn\_mpnn & -0.005 & [-0.007, -0.002] & 0.0008 & 0.010\\
task2 / smooth\_four\_genus / rf\_oracle vs rf\_diagram & -0.250 & [-0.265, -0.235] & 0.0002 & 0.010\\
task2 / smooth\_four\_genus / rf\_diagram vs linear\_cw & -0.138 & [-0.156, -0.119] & 0.0002 & 0.010\\
task2 / three\_genus / gnn\_mpnn vs rf\_diagram & -0.029 & [-0.043, -0.015] & 0.0002 & 0.010\\
task2 / three\_genus / gnn\_gin vs rf\_diagram & -0.086 & [-0.101, -0.070] & 0.0002 & 0.010\\
task2 / three\_genus / gnn\_gps vs rf\_diagram & -0.002 & [-0.015, +0.013] & 0.8184 & 1.000\\
task2 / three\_genus / gnn\_subgraph vs rf\_diagram & -0.079 & [-0.093, -0.066] & 0.0002 & 0.010\\
task2 / three\_genus / gnn\_mpnn vs linear\_cw & +0.095 & [+0.075, +0.116] & 0.0002 & 0.010\\
task2 / three\_genus / gnn\_gin vs linear\_cw & +0.039 & [+0.020, +0.058] & 0.0002 & 0.010\\
task2 / three\_genus / gnn\_gps vs linear\_cw & +0.123 & [+0.102, +0.143] & 0.0002 & 0.010\\
task2 / three\_genus / gnn\_subgraph vs linear\_cw & +0.045 & [+0.026, +0.065] & 0.0002 & 0.010\\
task2 / three\_genus / gnn\_gin vs gnn\_mpnn & -0.057 & [-0.062, -0.051] & 0.0002 & 0.010\\
task2 / three\_genus / rf\_oracle vs rf\_diagram & -0.687 & [-0.709, -0.666] & 0.0002 & 0.010\\
task2 / three\_genus / rf\_diagram vs linear\_cw & +0.124 & [+0.105, +0.144] & 0.0002 & 0.010\\
task2 / signature / gnn\_mpnn vs rf\_diagram & -0.022 & [-0.053, +0.008] & 0.1488 & 1.000\\
task2 / signature / gnn\_gin vs rf\_diagram & -0.059 & [-0.090, -0.030] & 0.0002 & 0.010\\
task2 / signature / gnn\_gps vs rf\_diagram & -0.006 & [-0.035, +0.024] & 0.7384 & 1.000\\
task2 / signature / gnn\_subgraph vs rf\_diagram & -0.058 & [-0.088, -0.030] & 0.0002 & 0.010\\
task2 / signature / gnn\_mpnn vs linear\_cw & +0.043 & [+0.019, +0.068] & 0.0012 & 0.010\\
task2 / signature / gnn\_gin vs linear\_cw & +0.006 & [-0.020, +0.032] & 0.6304 & 1.000\\
task2 / signature / gnn\_gps vs linear\_cw & +0.060 & [+0.034, +0.086] & 0.0002 & 0.010\\
task2 / signature / gnn\_subgraph vs linear\_cw & +0.007 & [-0.016, +0.031] & 0.5432 & 1.000\\
task2 / signature / gnn\_gin vs gnn\_mpnn & -0.037 & [-0.054, -0.020] & 0.0002 & 0.010\\
task2 / signature / rf\_oracle vs rf\_diagram & -0.134 & [-0.162, -0.108] & 0.0002 & 0.010\\
task2 / signature / rf\_diagram vs linear\_cw & +0.066 & [+0.040, +0.092] & 0.0002 & 0.010\\
\bottomrule
\end{longtable}
}
\section{Secondary Protocol (train $c{\le}11$, no validation)}
Reported in \texttt{results/rigorous\_tasks\_benchmark.json}
(\texttt{task2\_invariants}); superseded by Table~\ref{tab:task2}.
\section{Hyperparameters and Selection Rules}\todo{}
\section{Compute}\todo{hardware, CPU-hours, per-run time}
\section{Datasheet for KnotBench}\todo{motivation, composition, collection,
uses, distribution, maintenance; Croissant metadata}

\clearpage
\section*{NeurIPS Paper Checklist}
\todo{official checklist}

\end{document}
